This section lists and introduces the blogs I have posted or prepared
for the \href{https://www.facebook.com/groups/awsstudygroupfcj}{AWS
Study Group}. The blog entries are organized into three small sections,
following the same structure as the workshop sample.

\vspace{0.5em}\noindent\textbf{Blog 1 - WHY DOES AWS STILL CHARGE AFTER STOPPING
EC2?}\par

This blog explains why AWS may still generate charges even after an EC2
instance has been stopped, focusing on related billable resources that
may continue to exist outside the running instance itself.

\textbf{Published on:} July 22, 2026\\
\textbf{Platform:} AWS Study Groups\\
\textbf{Status:} Approved\\
\textbf{Public link:}
\href{https://www.facebook.com/groups/awsstudygroupfcj/permalink/2219932025438424/?rdid=SGRgenG5JS3UR6Y0\#}{Blog
1}

\vspace{0.5em}\noindent\textbf{Blog 2 - BUILDING A SCALABLE REAL-TIME CHAT SYSTEM WITH
SOCKET.IO, REDIS, AND
DYNAMODB}\par

This blog introduces the design of a scalable real-time chat system
using Socket.IO, Redis, and DynamoDB. It focuses on how real-time
communication, distributed state, and persistent storage can work
together in a larger system.

\textbf{Published on:} July 23, 2026\\
\textbf{Platform:} AWS Study Groups\\
\textbf{Status:} Pending

\vspace{0.5em}\noindent\textbf{Blog 3 - A Few Takeaways from My Software Engineering
Journey: From Writing Code to Breaking the System and Fixing It
Myself}\par

This article reflects on the journey from being someone who mainly
focused on writing code to becoming a Software Engineer with a broader
understanding of how systems operate, fail, and recover. Through working
with backend, frontend, Docker, Kubernetes, CI/CD, IAM, databases,
message queues, and AWS, the author realized that code running
successfully in a local environment is not enough. A stable system also
depends on configuration, networking, permissions, data, infrastructure,
and failure-handling mechanisms. The main lesson is to debug
systematically, identify the actual point of failure instead of making
random changes, and design systems so that the same problems are less
likely to happen again.

\textbf{Published date:} July 28, 2026\\
\textbf{Platform:} AWS Study Groups\\
\textbf{Status:} Pending\\
\textbf{Public link:}
\href{https://www.facebook.com/groups/awsstudygroupfcj/permalink/2227122848052675/\#}{Blog
3}

\vspace{0.5em}\noindent\textbf{Blog 4 - AWS Networking Advanced: 3 Techniques for
Building an Enterprise-Grade
VPC}\par

This article presents three techniques for upgrading a basic VPC
architecture into a design that is more suitable for enterprise systems.
It focuses on using VPC Endpoints to reduce traffic through NAT
Gateways, deploying infrastructure across multiple Availability Zones to
improve availability and reduce unnecessary cross-AZ data transfer, and
using AWS Transit Gateway to manage connections between multiple VPCs
through a centralized network model. The article demonstrates how
network architecture optimization can improve security, scalability,
operational control, and AWS cost efficiency.

\textbf{Published date:} July 29, 2026\\
\textbf{Platform:} AWS Study Groups\\
\textbf{Status:} Published\\
\textbf{Public link:}
\href{https://www.facebook.com/groups/awsstudygroupfcj/permalink/2225997548165205/\#}{Blog
4}
